% from ega4-2-fr.tex
\oldpage[IV-2]{58}

\subsection{Rappels sur les produits tensoriels de corps}
\label{subsection:IV.2.4.3-fr}

Pour la commodité du lecteur, nous rappelons ici certaines propriétés des
produits tensoriels de corps dont nous ferons usage dans les numéros suivants~;
pour les démonstrations, nous renvoyons à [1].

\begin{env}[4.3.1]
\label{IV.2.4.3.1-fr}
Rappelons qu'une extension $L$ d'un corps $k$ est dite \emph{primaire} si la
plus grande extension algébrique séparable de $k$ contenue dans $L$ est
$k$ lui-même.
\end{env}

\begin{proposition}[4.3.2]
\label{IV.2.4.3.2-fr}
Soient $K$, $L$ deux extensions d'un corps $k$. Si $L$ est une extension
primaire de $k$, alors $\operatorname{Spec}(L\otimes_k K)$ est irréductible,
et si $\xi$ est son point générique, $k(\xi)$ est une extension primaire de
$K$. Réciproquement, si pour toute extension finie séparable $K$ de $k$,
$\operatorname{Spec}(L\otimes_k K)$ est irréductible, $L$ est une extension
primaire de $k$.
\end{proposition}

\begin{proof}
Voir la démonstration dans [1], p.~14-03 à 14-06.
\end{proof}

\begin{corollary}[4.3.3]
\label{IV.2.4.3.3-fr}
Si $k$ est séparablement clos (i.e. si sa clôture algébrique est radicielle sur
$k$), alors, pour deux extensions quelconques $K$, $L$ de $k$,
$\operatorname{Spec}(L\otimes_k K)$ est irréductible, et réciproquement.
\end{corollary}

\begin{proof}
Cela découle aussitôt de (4.3.2).
\end{proof}

\begin{corollary}[4.3.4]
\label{IV.2.4.3.4-fr}
Soient $L$ une extension d'un corps $k$, $L_s$ la plus grande extension
algébrique séparable de $k$ contenue dans $L$ (parfois appelée la
\emph{fermeture algébrique séparable} de $k$ dans $L$)~; supposons $L_s$ de
degré fini sur $k$, et soit $K$ une extension galoisienne de $k$ contenant
$L_s$. Alors $\operatorname{Spec}(L\otimes_k K)$ a $[L_s:k]$ composantes
irréductibles dont il est la somme et les corps résiduels aux points génériques
de ces composantes sont des extensions primaires de $K$.
\end{corollary}

\begin{proof}
En effet, on a $L\otimes_k K=L\otimes_{L_s}(L_s\otimes_k K)$ et on sait que
$L_s\otimes_k K$ est isomorphe à un produit de $[L_s:k]$ corps isomorphes à
$K$ (Bourbaki, \emph{Alg.}, chap.~VIII, \S~7, n\textsuperscript{o}~3,
cor.~2 du th.~1)~; donc $L\otimes_k K$ est isomorphe à un produit de
$[L_s:k]$ anneaux isomorphes à $L\otimes_{L_s}K$~; comme $L$ est une
extension primaire de $L_s$, la conclusion résulte de (4.3.2).
\end{proof}

\begin{proposition}[4.3.5]
\label{IV.2.4.3.5-fr}
Soient $K$, $L$ deux extensions d'un corps $k$. Si $L$ est une extension
séparable de $k$, l'anneau $L\otimes_k K$ est réduit. Réciproquement, si pour
toute extension finie radicielle $K$ de $k$, l'anneau $L\otimes_k K$ est
réduit, $L$ est une extension séparable de $k$.
\end{proposition}

\begin{proof}
Pour la démonstration, voir Bourbaki, \emph{Alg.}, chap.~VIII, \S~7,
n\textsuperscript{o}~3, th.~1.
\end{proof}

\begin{corollary}[4.3.6]
\label{IV.2.4.3.6-fr}
Si $k$ est parfait, alors, pour deux extensions quelconques $K$, $L$ de $k$,
$K\otimes_k L$ est réduit, et réciproquement.
\end{corollary}

\begin{proof}
Cela résulte aussitôt de (4.3.5).
\end{proof}

\begin{corollary}[4.3.7]
\label{IV.2.4.3.7-fr}
Soit $L$ une extension séparable de $k$ et soit $K$ une extension de $k$ telle
que $K$ ou $L$ soit finie sur $k$, alors les corps résiduels de l'anneau
semi-local $L\otimes_k K$ sont des extensions séparables de $K$.
\end{corollary}

\begin{proof}
On sait en effet que $L\otimes_k K$ est composé direct de ces corps résiduels
$E_i$~; pour toute extension $K'$ de $K$, $L\otimes_k K'$ est donc isomorphe
au composé direct des anneaux $E_i\otimes_K K'$~; comme $L\otimes_k K'$ est
réduit, les $E_i$ sont séparables sur $K$ en vertu de (4.3.5).
\end{proof}

\begin{corollary}[4.3.8]
\label{IV.2.4.3.8-fr}
Si $K$ et $L$ sont deux extensions finies séparables de $k$, $K\otimes_k L$
est produit d'un nombre fini d'extensions séparables de $k$.
\end{corollary}

